% GENERATED FILE. Edit canonical lessons, then run npm run generate:books.
% canonical-source-hash: fnv1a64:d21caa117f6f7f33
% canonical-lessons: MW-C260-bahu, MW-C260-javai, MW-C260-bhavaj, MW-C260-bahnoi, MW-C260-salo

\chapter{Married Into the Family}
\label{ch:mw-married-into-the-family}
\hlchaptermodality{260}

\begin{chapteropening}
\textbf{By the end of this chapter:} \emph{I can name the in-laws who married into my family.}

Complete the last lesson of chapter 260: I can name the in-laws who married into my family.
\end{chapteropening}

\section[\texorpdfstring{\mw{बहू} (bahū)}{बहू (bahū)}]{\mw{बहू} (bahū) — a daughter-in-law; a young wife of the house}

\label{lesson:MW-C260-bahu}



\begin{quote}
Before the new one: say the Marwadi for a husband's sister, then the Marwadi for a husband (in his wife's words).
\end{quote}

\subsection*{You'll want to know: \mw{बहू}}
\textbf{\mw{बहू}} — \emph{bahū} — \textquotedblleft{}a daughter-in-law; a young wife of the house\textquotedblright{}.

\begin{quote}
\mw{म्हारी} \mw{बहू} \mw{घणी} \mw{हुसियार} \mw{है।}
\end{quote}

\begin{quote}
\emph{mhārī bahū ghaṇī husiyār hai.}
\end{quote}

\begin{quote}
My daughter-in-law is very clever.
\end{quote}

A feminine noun.

\begin{cousinweb}
\textbf{\mw{बहू}} (\emph{bahū}) is inherited from Sanskrit \emph{vadhū}, \textquotedblleft{}a bride, a young wife\textquotedblright{}. CDIAL gives no Marwari form in this entry; it lists Hindi \emph{bahū}, \textquotedblleft{}son's wife\textquotedblright{} (CDIAL 11250).
\end{cousinweb}

\begin{grammarlens}[title={What you've built}]
One more word for this chapter.
\end{grammarlens}

\subsection*{Your turn}
\begin{itemize}
\raggedright
  \item \emph{Say it:} \emph{bahū}
  \item \emph{Say it:} \emph{bahū}, once more
  \item \emph{Say it:} \emph{mhārī bahū ghaṇī husiyār hai.} — My daughter-in-law is very clever.
  \item \emph{Recall:} say the Marwadi for a husband's sister, then the Marwadi for a husband (in his wife's words), then say \emph{bahū} again
\end{itemize}

\begin{tcolorbox}[breakable,colback=teal!4,colframe=teal!35!black,title={Before you move on}]
What does \mw{बहू} mean? (\textquotedblleft{}A daughter-in-law; a young wife of the house\textquotedblright{}.) Say it once more.
\end{tcolorbox}
\section[\texorpdfstring{\mw{जंवाई} (jãvāī)}{जंवाई (jãvāī)}]{\mw{जंवाई} (jãvāī) — a son-in-law}

\label{lesson:MW-C260-javai}



\begin{quote}
Before the new one: say the Marwadi for a husband (in his wife's words), then the Marwadi for a daughter-in-law.
\end{quote}

\subsection*{You'll want to know: \mw{जंवाई}}
\textbf{\mw{जंवाई}} — \emph{jãvāī} — \textquotedblleft{}a son-in-law\textquotedblright{}.

\begin{quote}
\mw{जंवाई} \mw{काल} \mw{आयो।}
\end{quote}

\begin{quote}
\emph{jãvāī kāl āyo.}
\end{quote}

\begin{quote}
Our son-in-law came yesterday.
\end{quote}

A masculine noun.

\begin{cousinweb}
\textbf{\mw{जंवाई}} (\emph{jãvāī}) is inherited from Sanskrit \emph{jāmātṛ}, \textquotedblleft{}a daughter's husband\textquotedblright{}. CDIAL gives no Marwari form in this entry; it lists Hindi \emph{jãwāī} (CDIAL 5198).
\end{cousinweb}

\begin{grammarlens}[title={What you've built}]
One more word for this chapter.
\end{grammarlens}

\subsection*{Your turn}
\begin{itemize}
\raggedright
  \item \emph{Say it:} \emph{jãvāī}
  \item \emph{Say it:} \emph{jãvāī}, once more
  \item \emph{Say it:} \emph{jãvāī kāl āyo.} — Our son-in-law came yesterday.
  \item \emph{Recall:} say the Marwadi for a husband (in his wife's words), then the Marwadi for a daughter-in-law, then say \emph{jãvāī} again
\end{itemize}

\begin{tcolorbox}[breakable,colback=teal!4,colframe=teal!35!black,title={Before you move on}]
What does \mw{जंवाई} mean? (\textquotedblleft{}A son-in-law\textquotedblright{}.) Say it once more.
\end{tcolorbox}
\section[\texorpdfstring{\mw{भावज} (bhāvaj)}{भावज (bhāvaj)}]{\mw{भावज} (bhāvaj) — a brother's wife}

\label{lesson:MW-C260-bhavaj}



\begin{quote}
Before the new one: say the Marwadi for a daughter-in-law, then the Marwadi for a son-in-law.
\end{quote}

\subsection*{You'll want to know: \mw{भावज}}
\textbf{\mw{भावज}} — \emph{bhāvaj} — \textquotedblleft{}a brother's wife\textquotedblright{}.

\begin{quote}
\mw{म्हारी} \mw{भावज} \mw{रोटी} \mw{पकावै} \mw{है।}
\end{quote}

\begin{quote}
\emph{mhārī bhāvaj roṭī pakāvai hai.}
\end{quote}

\begin{quote}
My brother's wife is cooking flatbread.
\end{quote}

A feminine noun.

\begin{cousinweb}
\textbf{\mw{भावज}} (\emph{bhāvaj}) is inherited from Sanskrit \emph{bhrātṛjāyā}, \textquotedblleft{}a brother's wife\textquotedblright{}; CDIAL lists Old Marwari \emph{bhāvaja} under it (CDIAL 9660).
\end{cousinweb}

\begin{grammarlens}[title={What you've built}]
One more word for this chapter.
\end{grammarlens}

\subsection*{Your turn}
\begin{itemize}
\raggedright
  \item \emph{Say it:} \emph{bhāvaj}
  \item \emph{Say it:} \emph{bhāvaj}, once more
  \item \emph{Say it:} \emph{mhārī bhāvaj roṭī pakāvai hai.} — My brother's wife is cooking flatbread.
  \item \emph{Recall:} say the Marwadi for a daughter-in-law, then the Marwadi for a son-in-law, then say \emph{bhāvaj} again
\end{itemize}

\begin{tcolorbox}[breakable,colback=teal!4,colframe=teal!35!black,title={Before you move on}]
What does \mw{भावज} mean? (\textquotedblleft{}A brother's wife\textquotedblright{}.) Say it once more.
\end{tcolorbox}
\section[\texorpdfstring{\mw{बहनोई} (bahnoī)}{बहनोई (bahnoī)}]{\mw{बहनोई} (bahnoī) — a sister's husband}

\label{lesson:MW-C260-bahnoi}



\begin{quote}
Before the new one: say the Marwadi for a son-in-law, then the Marwadi for a brother's wife.
\end{quote}

\subsection*{You'll want to know: \mw{बहनोई}}
\textbf{\mw{बहनोई}} — \emph{bahnoī} — \textquotedblleft{}a sister's husband\textquotedblright{}.

\begin{quote}
\mw{म्हारो} \mw{बहनोई} \mw{मास्टर} \mw{है।}
\end{quote}

\begin{quote}
\emph{mhāro bahnoī māsṭar hai.}
\end{quote}

\begin{quote}
My sister's husband is a teacher.
\end{quote}

A masculine noun.

\begin{cousinweb}
\textbf{\mw{बहनोई}} (\emph{bahnoī}) is inherited from Sanskrit \emph{bhaginīpati}, \textquotedblleft{}a sister's husband\textquotedblright{}; CDIAL lists Old Marwari \emph{bahaṇeu} under it (CDIAL 9350).
\end{cousinweb}

\begin{grammarlens}[title={What you've built}]
One more word for this chapter.
\end{grammarlens}

\subsection*{Your turn}
\begin{itemize}
\raggedright
  \item \emph{Say it:} \emph{bahnoī}
  \item \emph{Say it:} \emph{bahnoī}, once more
  \item \emph{Say it:} \emph{mhāro bahnoī māsṭar hai.} — My sister's husband is a teacher.
  \item \emph{Recall:} say the Marwadi for a son-in-law, then the Marwadi for a brother's wife, then say \emph{bahnoī} again
\end{itemize}

\begin{tcolorbox}[breakable,colback=teal!4,colframe=teal!35!black,title={Before you move on}]
What does \mw{बहनोई} mean? (\textquotedblleft{}A sister's husband\textquotedblright{}.) Say it once more.
\end{tcolorbox}
\section[\texorpdfstring{\mw{साळो} (sāḷo)}{साळो (sāḷo)}]{\mw{साळो} (sāḷo) — a wife's brother}

\label{lesson:MW-C260-salo}



\begin{quote}
Before the new one: say the Marwadi for a brother's wife, then the Marwadi for a sister's husband.
\end{quote}

\subsection*{You'll want to know: \mw{साळो}}
\textbf{\mw{साळो}} — \emph{sāḷo} — \textquotedblleft{}a wife's brother\textquotedblright{}.

\begin{quote}
\mw{म्हारो} \mw{साळो} \mw{पुलिस} \mw{मांय} \mw{है।}
\end{quote}

\begin{quote}
\emph{mhāro sāḷo pulis mā̃y hai.}
\end{quote}

\begin{quote}
My wife's brother is in the police.
\end{quote}

A masculine noun.

\begin{cousinweb}
\textbf{\mw{साळो}} (\emph{sāḷo}) is inherited from Sanskrit \emph{syāla}, \textquotedblleft{}a wife's brother\textquotedblright{}. CDIAL gives no Marwari form in this entry; it lists Gujarati \emph{sāḷo}, with this book's \mw{ळ} (CDIAL 13871).
\end{cousinweb}

\begin{grammarlens}[title={What you've built}]
One more word for this chapter.
\end{grammarlens}

\subsection*{Your turn}
\begin{itemize}
\raggedright
  \item \emph{Say it:} \emph{sāḷo}
  \item \emph{Say it:} \emph{sāḷo}, once more
  \item \emph{Say it:} \emph{mhāro sāḷo pulis mā̃y hai.} — My wife's brother is in the police.
  \item \emph{Recall:} say the Marwadi for a brother's wife, then the Marwadi for a sister's husband, then say \emph{sāḷo} again
\end{itemize}

\begin{tcolorbox}[breakable,colback=teal!4,colframe=teal!35!black,title={Before you move on}]
What does \mw{साळो} mean? (\textquotedblleft{}A wife's brother\textquotedblright{}.) Say it once more.
\end{tcolorbox}
